\documentclass[11pt]{article}
\usepackage[T5]{fontenc}
\usepackage[margin=2.5cm]{geometry}
\usepackage{graphicx}

\title{Labwork 2: GPU Device Info}
\author{Phạm Đỗ Ngọc Minh -- 2540075}
\date{}

\begin{document}
\maketitle

My laptop is a MacBook Pro with an Apple M1 chip. Its GPU is not an NVIDIA GPU,
so \texttt{numba.cuda} cannot be used. I followed the same steps with Apple's
Metal API instead (\texttt{gpu\_info\_mac.py}): list the GPUs, select the default
one, then read its id, name, core info and memory info. Metal does not report
the core count, so it is read from the IORegistry (\texttt{ioreg}); one Apple GPU
core plays the role of a CUDA multiprocessor and contains 128 ALUs.

\begin{figure}[h]
  \centering
  \includegraphics[width=0.85\textwidth]{../gpuinformation.png}
  \caption{GPU information of the Apple M1}
\end{figure}

\end{document}
